\documentclass[12pt]{article}
\usepackage[utf8]{inputenc}
\usepackage[T1]{fontenc}
\usepackage{hyperref}
\title{Tembo Language Resources for K-12 Curriculum Development}
\author{Compiled by Claude Sonnet 5 (claude-sonnet-5) for Andr\'as Kornai}
\date{}
\begin{document}
\maketitle

\section{Introduction}
Tembo (Kitembo, Chitembo; ISO 639-3 tbt, Glottocode temb1270, Guthrie JD.531) is a Bantu language of the Kivu provinces of eastern DR Congo. The speakers are called Batembo. Sources place it mainly in Masisi and Walikale (North Kivu) and Kalehe (South Kivu). Speaker estimates conflict: about 150,000 (1994) against 500,000 to 700,000 (2023, per Wikipedia). Search results mention a Latin alphabet devised by Belgian missionaries, mainly for Bible translation, and a separate Hangul-based alphabet for the Twa community (a Korean-linked initiative from 2015).

\section{Primary Educational Resources}
No primers, textbooks or school materials were found in search results.

\section{Linguistic and Reference Materials}
\begin{itemize}
\item Wikipedia, Tembo (Kitembo) language: \url{https://en.wikipedia.org/wiki/Tembo_(Kitembo)_language}
\item Omniglot, Tembo (Kitembo / Chitembo) (writing system notes): \url{https://www.omniglot.com/writing/tembo.htm}
\item Glottolog, Tembo: \url{https://glottolog.org/resource/languoid/id/temb1270.bigmap.html}
\end{itemize}

\section{Religious and Community Resources}
Search results indicate that the New Testament has a 2017 edition. I did not find a direct link to it.

\section{Digital and Audio Resources}
\begin{itemize}
\item Global Recordings Network, Kitembo (audio and video for evangelism and Bible teaching): \url{https://globalrecordings.net/en/language/tbt}
\end{itemize}

\section{Curriculum Development Context}
Search found no evidence of Tembo in schools. A Latin orthography exists, which helps. No pedagogical materials were found. French, Swahili and other national languages are the likely school languages in the area.
\end{document}